\chapter{Local Iwasawa Selmer}
\label{chap:local-iwasawa-selmer}

% archon:covers MR4372220OnTheAnticyclotomicIwasawaTheoryOfRationalEllipticCurvesAtEisensteinPrimes/Local.lean

This chapter fixes the local anticyclotomic Iwasawa notation once, packages the character Selmer groups as kernels of localisation maps, and isolates the split and \(p\)-adic local cohomology input needed later by the algebraic and Howard branches.

\begin{definition}[Iwasawa algebra and Selmer structures]
  \label{def:iwasawa-algebra}
  \lean{MR4372220.Local.iwasawaAlgebra}
  % SOURCE: [Castella--Grossi--Lee--Skinner, Local cohomology groups of characters; Selmer groups of characters; Local cohomology groups of E; Selmer groups of E] (read from references/MR4372220-anticyclotomic-iwasawa-source/Eisenstein.tex)
  % SOURCE QUOTE: "Set \(M_\theta=\Z_p(\theta)\otimes_{\Z_p}\Lambda^\vee\),"
  \textit{Source: Castella--Grossi--Lee--Skinner, Local cohomology groups of characters; Selmer groups of characters; Local cohomology groups of E; Selmer groups of E.}
  Let \(\Gamma = \mathrm{Gal}(K_\infty/K)\) and \(\Lambda = \mathbb{Z}_p\llbracket \Gamma\rrbracket\), and write \(\Lambda^{\mathrm{ur}} = \Lambda \widehat{\otimes}_{\mathbb{Z}_p}\mathbb{Z}_p^{\mathrm{ur}}\).
  For a character \(\theta : G_K \to \mathbb{F}_p^\times\), set
  \[
    M_\theta = \mathbb{Z}_p(\theta)\otimes_{\mathbb{Z}_p}\Lambda^\vee,
  \]
  where \(\Lambda^\vee = \operatorname{Hom}_{\mathrm{cts}}(\Lambda,\mathbb{Q}_p/\mathbb{Z}_p)\), and for an elliptic curve \(E/\mathbb{Q}\) of good reduction at \(p\), set
  \[
    M_E = T_pE\otimes_{\mathbb{Z}_p}\Lambda^\vee.
  \]
  For a finite set \(\Sigma\) of places of \(K\) containing \(\infty\), the places above \(p\), the primes dividing the relevant conductor, and all finite places split in \(K\), define the Greenberg Selmer group by the kernel of the usual global-to-local map
  \[
    H^1_{\mathcal{F}_{\mathrm{Gr}}}(K,M_\theta)
    = \ker\!\left(H^1(K^\Sigma/K,M_\theta)\to \prod_{w\in\Sigma,\;w\nmid p} H^1(K_w,M_\theta)\times H^1(K_{\bar v},M_\theta)\right),
  \]
  and the \(S\)-imprimitive variant, where \(S=\Sigma\setminus\{v,\bar v,\infty\}\), by
  \[
    H^1_{\mathcal{F}_{\mathrm{Gr}}^S}(K,M_\theta)
    = \ker\!\left(H^1(K^\Sigma/K,M_\theta)\to H^1(K_{\bar v},M_\theta)\right).
  \]
  Replacing \(M_\theta\) with \(M_\theta[p]\) gives the residual character Selmer groups \(H^1_{\mathcal{F}_{\mathrm{Gr}}}(K,M_\theta[p])\) and \(H^1_{\mathcal{F}_{\mathrm{Gr}}^S}(K,M_\theta[p])\). Their Pontryagin duals are denoted \(\mathfrak{X}_\theta\) and \(\mathfrak{X}_\theta^S\). The corresponding primitive and \(S\)-imprimitive Selmer groups for \(M_E\) are defined by the same kernel conditions, but this chapter does not enter the \(M_E[p]\) comparison theory.
\end{definition}

\begin{definition}[Characteristic ideal]
  \label{def:char-ideal}
  \lean{MR4372220.Local.charIdeal}
  \uses{def:iwasawa-algebra}
  % SOURCE: [Castella--Grossi--Lee--Skinner, general commutative algebra for Iwasawa modules] (read from references/MR4372220-anticyclotomic-iwasawa-source/Eisenstein.tex)
  % SOURCE QUOTE: "char_\Lambda(\mathfrak{X}_E)"
  \textit{Source: Standard Iwasawa theory, characteristic ideals of torsion \(\Lambda\)-modules.}

  Let \(M\) be a finitely generated torsion \(\Lambda\)-module. The \emph{characteristic ideal} of \(M\) is
  \[
    \mathrm{char}_\Lambda(M) = \prod_{\mathfrak{P}} \mathfrak{P}^{\mathrm{length}_{\Lambda_\mathfrak{P}}(M_\mathfrak{P})} \subset \Lambda,
  \]
  where the product runs over height-one primes \(\mathfrak{P}\) of \(\Lambda\). It is a principal ideal \(\mathrm{char}_\Lambda(M) = (f_M)\) for some \(f_M \in \Lambda\) (since \(\Lambda\) is a UFD). The \emph{\(\mu\)-invariant} of \(M\) is \(\mu(M) = \mathrm{ord}_{(p)}(f_M)\) and the \emph{\(\lambda\)-invariant} is \(\lambda(M) = \deg(\bar{f}_M)\) where \(\bar{f}_M \in \mathbb{Z}_p[T]\) is the image of \(f_M\) after the choice of topological generator \(\gamma \in \Gamma\).

  Key properties: (1) \(\mathrm{char}_\Lambda\) is multiplicative on short exact sequences of torsion \(\Lambda\)-modules; (2) \(\mu(M) = 0\) and \(\lambda(M) = n\) iff \(\mathrm{length}_{\mathbb{F}_p}(M[p]) = n\) (when \(M\) has no non-trivial finite submodules); (3) if \(M\) and \(N\) are pseudo-isomorphic then \(\mathrm{char}_\Lambda(M) = \mathrm{char}_\Lambda(N)\).
\end{definition}

\begin{lemma}[Commutative algebra criterion]
  \label{lem:commalg}
  \lean{MR4372220.Local.commAlgCriterion}
  \uses{def:iwasawa-algebra}
  % SOURCE: [Castella--Grossi--Lee--Skinner, Local cohomology groups of characters, \(w\mid p\)] (read from references/MR4372220-anticyclotomic-iwasawa-source/Eisenstein.tex)
  % SOURCE QUOTE: "a free \(\Lambda\)-module"
  \textit{Source: Castella--Grossi--Lee--Skinner, Local cohomology groups of characters, \(w\mid p\).}
  Let \(X\) be a finitely generated \(\Lambda\)-module. If \(X[T]=0\) and \(X/TX\) is a free \(\mathbb{Z}_p\)-module of rank \(r\), then \(X\) is a free \(\Lambda\)-module of rank \(r\).
\end{lemma}

\begin{proof}
Use Nakayama to produce a surjection \(\Lambda^r\twoheadrightarrow X\), then compare modulo \(T\). The condition \(X[T]=0\) kills the kernel after reduction, and a second Nakayama argument shows that the kernel is already zero.
\end{proof}

For \(w\nmid p\) split in \(K\), write \(\mathcal{P}_w(\theta)\in\Lambda\) for the Euler factor attached to \(w\).

\begin{lemma}[Split local cohomology away from \(p\)]
  \label{lem:local-char-wneqp}
  \lean{MR4372220.Local.localCohomologySplit}
  \uses{def:iwasawa-algebra}
  % SOURCE: [Castella--Grossi--Lee--Skinner, Local cohomology groups of characters, \(w\nmid p\) split in \(K\)] (read from references/MR4372220-anticyclotomic-iwasawa-source/Eisenstein.tex)
  % SOURCE QUOTE: "char_\Lambda(\rH^1(K_w,M_\theta)^\vee)=(\mathcal{P}_w(\theta))"
  \textit{Source: Castella--Grossi--Lee--Skinner, Local cohomology groups of characters, \(w\nmid p\) split in \(K\).}
  For a prime \(w\nmid p\) split in \(K\), the local dual \(H^1(K_w,M_\theta)^\vee\) is \(\Lambda\)-torsion with characteristic ideal generated by \(\mathcal{P}_w(\theta)\). In particular, \(H^1(K_w,M_\theta)^\vee\) has \(\mu=0\).
\end{lemma}

\begin{proof}
The key point is that \(w\) is finitely decomposed in \(K_\infty/K\), so the local cohomology is governed by the unramified quotient of inertia. The characteristic ideal is the local Euler factor, and therefore no nontrivial \(\mu\)-part can appear.
\end{proof}

\begin{proposition}[Local cohomology at \(p\)]
  \label{lem:local-char-wp}
  \lean{MR4372220.Local.localCohomologyAtP}
  \uses{def:iwasawa-algebra, lem:commalg}
  % SOURCE: [Castella--Grossi--Lee--Skinner, Local cohomology groups of characters, \(w\mid p\)] (read from references/MR4372220-anticyclotomic-iwasawa-source/Eisenstein.tex)
  % SOURCE QUOTE: "\(\Lambda\)-cofree of rank \(1\)"
  \textit{Source: Castella--Grossi--Lee--Skinner, Local cohomology groups of characters, \(w\mid p\).}
  Let \(\omega : G_{\mathbb{Q}}\to \mathbb{F}_p^\times\) be the mod \(p\) cyclotomic character, and let \(w\mid p\) be a place of \(K\). If \(\theta|_{G_w}\neq 1,\omega\), then the restriction map
  \[
    H^1(K_w,M_\theta)\to H^1(I_w,M_\theta)^{G_w/I_w}
  \]
  is an isomorphism, and \(H^1(K_w,M_\theta)\) is \(\Lambda\)-cofree of rank \(1\).
\end{proposition}

\begin{proof}
Injectivity comes from the vanishing of \(M_\theta^{G_w}\), and the cofree statement is reduced to a rank computation for the Pontryagin dual. The quotient by \(T\) vanishes on the character module, so \cref{lem:commalg} applies after local duality and the local Euler characteristic formula.
\end{proof}

\begin{theorem}[Rubin--Hida \(\mu=0\) theorem]
  \label{thm:mu-zero-theta}
  \lean{MR4372220.Local.muZeroTheta}
  \uses{def:iwasawa-algebra, lem:local-char-wneqp, lem:local-char-wp}
  % SOURCE: [Castella--Grossi--Lee--Skinner, Selmer groups of characters] (read from references/MR4372220-anticyclotomic-iwasawa-source/Eisenstein.tex)
  % SOURCE QUOTE: "torsion \(\Lambda\)-module with \(\mu\)-invariant zero"
  \textit{Source: Castella--Grossi--Lee--Skinner, Selmer groups of characters.}
  If \(\theta|_{G_{\bar v}}\neq 1,\omega\), then \(\mathfrak{X}_\theta\) is a torsion \(\Lambda\)-module with \(\mu=0\).
\end{theorem}

\begin{proof}
Identify the Selmer group with the \(\theta\)-isotypic component of the relevant anticyclotomic class group by Shapiro's lemma and class field theory. Rubin's main-conjecture input gives the characteristic ideal, and Hida's vanishing of the \(\mu\)-invariant for the Katz \(p\)-adic \(L\)-function gives the conclusion.
\end{proof}

\begin{proposition}[Imprimitive \(\lambda\)-formula]
  \label{prop:lambda-theta}
  \lean{MR4372220.Local.lambdaTheta}
  \uses{def:iwasawa-algebra, thm:mu-zero-theta, lem:local-char-wneqp, lem:local-char-wp}
  % SOURCE: [Castella--Grossi--Lee--Skinner, Selmer groups of characters] (read from references/MR4372220-anticyclotomic-iwasawa-source/Eisenstein.tex)
  % SOURCE QUOTE: "has no nonzero pseudo-null"
  \textit{Source: Castella--Grossi--Lee--Skinner, Selmer groups of characters.}
  If \(\theta|_{G_{\bar v}}\neq 1,\omega\), then \(\mathfrak{X}_\theta^S\) is a torsion \(\Lambda\)-module with \(\mu=0\), and
  \[
    \lambda(\mathfrak{X}_\theta^S)
    = \lambda(\mathfrak{X}_\theta) + \sum_{w\in S,\; w\nmid p}\lambda(\mathcal{P}_w(\theta)).
  \]
  Moreover, \(H^1_{\mathcal{F}_{\mathrm{Gr}}^S}(K,M_\theta[p])\) is finite and
  \[
    \dim_{\mathbb{F}_p}\bigl(H^1_{\mathcal{F}_{\mathrm{Gr}}^S}(K,M_\theta[p])\bigr)=\lambda(\mathfrak{X}_\theta^S).
  \]
\end{proposition}

\begin{proof}
Compare the primitive and \(S\)-imprimitive Selmer groups by exactness of localisation, then use \cref{thm:mu-zero-theta} and the split-prime local calculation to read off the Euler-factor contribution to \(\lambda\). The absence of nonzero pseudo-null submodules turns the \(\lambda\)-invariant into the size of the residual Selmer group.
\end{proof}

\begin{corollary}[Residual exactness]
  \label{cor:characters}
  \lean{MR4372220.Local.charactersResidualExactness}
  \uses{def:iwasawa-algebra, prop:lambda-theta, lem:local-char-wp}
  % SOURCE: [Castella--Grossi--Lee--Skinner, Corollary \ref{corcharacters}] (read from references/MR4372220-anticyclotomic-iwasawa-source/Eisenstein.tex)
  % SOURCE QUOTE: "is exact"
  \textit{Source: Castella--Grossi--Lee--Skinner, Corollary (Residual Exactness).}
  If \(\theta|_{G_{\bar v}}\neq 1,\omega\), then \(H^2(K^\Sigma/K,M_\theta[p])=0\), and the residual sequence
  \[
    0\to H^1_{\mathcal{F}_{\mathrm{Gr}}^S}(K,M_\theta[p])\to H^1(K^\Sigma/K,M_\theta[p])\to H^1(K_{\bar v},M_\theta[p])\to 0
  \]
  is exact.
\end{corollary}

\begin{proof}
The previous proposition gives \(H^2(K^\Sigma/K,M_\theta)=0\), and the short exact sequence induced by multiplication by \(p\) then identifies the residual \(H^2\)-group. Divisibility of the imprimitive Selmer group and the \(\Lambda\)-cofreeness of the local term at \(\bar v\) force the snake-lemma sequence to be exact.
\end{proof}
